% 06_results.tex — Results
\section{Results}
\label{sec:results}

\subsection{Experiment 1: Encoding Layer Validation}

The multimodal encoding layer achieves an average fidelity of \AVGFIDELITY{} across knowledge types (Figure~\ref{fig:fidelity}).
Policy encoding is lossless (\FIDELITYPOLICY{}), expert judgment encoding preserves all semantic content (\FIDELITYEXPERT{}), and quantitative encoding achieves \FIDELITYQUANT{} fidelity---the slight loss reflecting dimensionality reduction in distribution fitting.

Cross-type query performance strongly differentiates the encoding approach from baselines (Figure~\ref{fig:crosstype}).
The full system achieves \CROSSTYPEACCURACY{} accuracy on queries requiring reasoning across two or more knowledge types, compared to \RAWCONCATACCURACY{} for raw concatenation and \LLMBASELINEACCURACY{} for the LLM baseline.
The encoding layer's preservation of epistemic metadata enables the cross-encoder to correctly weigh conflicting signals---a capability that raw text concatenation cannot support.

The system detected \CONFLICTSDETECTED{}/\CONFLICTSTOTAL{} planted cross-source conflicts, including the price-promotion trap where expert judgment contradicts joint quantitative analysis, and the constraint-coupling conflict where expert recommendations exceed policy margin thresholds.

\begin{figure}[ht]
    \centering
    \includegraphics[width=0.85\textwidth]{figures/fig5_fidelity_heatmap.pdf}
    \caption{Encoding fidelity across knowledge types and encoding dimensions.}
    \label{fig:fidelity}
\end{figure}

\begin{figure}[ht]
    \centering
    \includegraphics[width=0.7\textwidth]{figures/fig6_cross_type_comparison.pdf}
    \caption{Cross-type query success rates comparing the full system against baselines.}
    \label{fig:crosstype}
\end{figure}

\subsection{Experiment 2: Ablation Study}

The ablation study demonstrates that each structural invariant contributes to system performance (Figure~\ref{fig:ablation}, Table~\ref{tab:ablation}).

The full system discovers \EFFECTSDISCOVERED{}/\EFFECTSTOTAL{} ground-truth effects with a precision of \FULLPRECISION{} and recall of \FULLRECALL{}.
Removing the coupling model (\textit{No Coupling}) reduces effects discovered to \NOCOUPLINGEFFECTS{}, as interaction effects GT-1 (price-promotion) and GT-2 (assortment-distribution) require cross-lever analysis.
Removing multimodal encoding (\textit{No Encoding}) reduces precision, as the system cannot properly weigh conflicting evidence without epistemic metadata.
Removing the progressive pipeline (\textit{No Pipeline}) yields the lowest precision at \NOPIPELINEPRECISION{}, confirming that unstructured single-pass analysis cannot systematically explore the decision space.

\begin{figure}[ht]
    \centering
    \includegraphics[width=\textwidth]{figures/fig4_ablation_bars.pdf}
    \caption{Ablation results: effects discovered, precision, and recall across four configurations.}
    \label{fig:ablation}
\end{figure}

\begin{figure}[ht]
    \centering
    \includegraphics[width=0.85\textwidth]{figures/fig9_discovery_matrix.pdf}
    \caption{Ground-truth discovery matrix: which effects are found by each configuration.}
    \label{fig:discovery}
\end{figure}

\begin{table}[ht]
\centering
\caption{Ablation study results across four system configurations.}
\label{tab:ablation}
\small
\begin{tabular}{lccccc}
\toprule
\textbf{Configuration} & \textbf{Effects (/5)} & \textbf{Precision} & \textbf{Recall} & \textbf{Violations} & \textbf{Revenue $\Delta$} \\
\midrule
Full System & 4 & 0.67 & 0.80 & 0 & -22% \\
No Coupling & 4 & 0.67 & 0.80 & 0 & -22% \\
No Encoding & 4 & 0.80 & 0.80 & 0 & -22% \\
No Pipeline & 4 & 0.40 & 0.80 & 0 & -22% \\
\bottomrule
\end{tabular}
\end{table}

\subsection{Experiment 3: Progressive Reduction Pipeline}

The pipeline achieves dramatic compression while preserving ground-truth coverage at every stage (Figure~\ref{fig:compression}).
Stage~1 (diagnostic agents) reduces the ${\sim}10^{18}$ input space to \PIPELINESTAGE{1} evidence packets.
Stage~2 (causal synthesis) consolidates to \PIPELINESTAGE{2} structured interventions.
Stage~3 (quality gate) selects \PIPELINEOUTPUT{} final recommendations with zero hard-constraint violations.

Ground-truth coverage remains at \GTCOVERAGE{3} through all stages, confirming that the compression is lossless with respect to the planted interaction effects.

The quality profile of the top interventions (Figure~\ref{fig:radar}) shows strong performance across all scoring dimensions, with perfect constraint alignment and high evidence density.

\begin{figure}[ht]
    \centering
    \includegraphics[width=0.65\textwidth]{figures/fig8_quality_radar.pdf}
    \caption{Quality gate radar chart for top-5 interventions across scoring dimensions.}
    \label{fig:radar}
\end{figure}

\begin{table}[ht]
\centering
\caption{Top recommended interventions with quality scores.}
\label{tab:interventions}
\small
\begin{tabular}{clcccl}
\toprule
\textbf{Rank} & \textbf{Lever} & \textbf{Quality} & \textbf{Evidence} & \textbf{Constraints} & \textbf{Revenue} \\
\midrule
1 & pricing & 0.91 & 1.00 & 1.00 & +28.1% \\
2 & pack\_price & 0.91 & 1.00 & 1.00 & +27.3% \\
3 & assortment & 0.90 & 1.00 & 1.00 & +27.2% \\
4 & distribution & 0.81 & 0.73 & 1.00 & +40.6% \\
5 & promotion & 0.73 & 1.00 & 1.00 & +74.8% \\
6 & unknown & 0.59 & 0.73 & 1.00 & +36.9% \\
\bottomrule
\end{tabular}
\end{table}

\subsection{Experiment 4: Cross-Domain Generalization}

The framework generalizes to two secondary domains using the same pipeline architecture (Figure~\ref{fig:generalization}, Table~\ref{tab:generalization}).
In a precision medicine mini-domain (3 coupled treatment levers, 100 patients), the system discovers 1/1 planted cross-lever interaction effects.
In a supply chain design mini-domain (3 coupled logistics levers, 50 nodes), the system again discovers 1/1 planted effects.

\begin{figure}[ht]
    \centering
    \includegraphics[width=\textwidth]{figures/fig10_generalization.pdf}
    \caption{Cross-domain generalization: effects discovered and pipeline functionality across three domains.}
    \label{fig:generalization}
\end{figure}

\begin{table}[ht]
\centering
\caption{Cross-domain generalization results.}
\label{tab:generalization}
\small
\begin{tabular}{lccc}
\toprule
\textbf{Domain} & \textbf{Effects Found} & \textbf{Effects Total} & \textbf{Pipeline Functional} \\
\midrule
RGM (BevCo) & 4 & 5 & \checkmark \\
Precision Medicine & 1 & 1 & \checkmark \\
Supply Chain & 1 & 1 & \checkmark \\
\bottomrule
\end{tabular}
\end{table}
